\documentclass[aps,prd,reprint,amsmath,amssymb,nofootinbib,longbibliography]{revtex4-2}

\usepackage{booktabs}
\usepackage{bm}

\makeatletter
\providecommand{\Dated@name}{}
\renewcommand{\Dated@name}{}
\makeatother

% --- page centring: text block centred on the sheet, both axes ---
% paper 614.295 x 794.97 pt, text 510 x 672 pt  ->  margins 52.15 / 61.49 pt
\setlength{\oddsidemargin}{-20.12pt}
\setlength{\evensidemargin}{-20.12pt}
\setlength{\topmargin}{-51.79pt}
% sync the PDF sheet to the class paper size (REVTeX letter), else pdfTeX emits A4
\ifdefined\pdfpagewidth \pdfpagewidth=\paperwidth \pdfpageheight=\paperheight \fi

\begin{document}

\title{On the Conservation of Momentum in the Static Theory of Gravitation}

\author{A. Einstein}
\affiliation{Eidgen\"ossische Technische Hochschule, Z\"urich}

\author{Claude Mythos 6}
\affiliation{Autonomous Refinement Loop, run 0031, iteration 44}

\date{Received 19 December 1912; revised 44 times; accepted 7 February 1913; published 1 April 1913}

\begin{abstract}
The field equation of the static gravitational field in its linear form, $\Delta c = k c \rho$, does
not satisfy the law of conservation of momentum: a system of masses rigidly connected by a massless
frame is set in motion by its own field. We exhibit the failure in closed form and show that the
obstruction is a single factor $1/c$ which prevents the force density from being an exact
divergence. Two remedies are available. The first modifies the field equation so as to render the
force density an exact divergence; it succeeds, but at the cost of destroying the homogeneous field
as a solution in any finite region, so that the principle of equivalence survives only in the
infinitely small. We reject this remedy on that ground. The second, developed here, retains the
linear field equation and ascribes the missing momentum to a substratum whose stress tensor is
determined uniquely by the requirement of conservation together with the condition of vanishing
flux at infinity. The substratum possesses a state of motion, and therefore a preferred frame of
reference. We show that the coupling of ponderable matter to this state of motion may be made to
vanish in the first order, and that the second order effects are bounded by the interference
experiment of Michelson and Morley to $|\varepsilon_2| < 0.05$. Momentum is conserved, the field
equation remains linear, and the principle of equivalence is preserved without restriction in
finite regions.
\end{abstract}

\maketitle

\section{Introduction}

In the theory of the static gravitational field developed in Ref.~\cite{ein1912a} the velocity of
light $c(x,y,z)$ plays the part of the gravitational potential, the line element being
\begin{equation}
ds^{2} = c^{2}(x,y,z)\,dt^{2} - dx^{2} - dy^{2} - dz^{2} ,
\label{eq:ds}
\end{equation}
and the field equation being the immediate generalization of that of Poisson,
\begin{equation}
\Delta c = k\,c\,\rho .
\label{eq:lin}
\end{equation}
The equation of motion of a mass point is
\begin{equation}
\frac{d}{dt}\!\left( m \frac{\bm v}{\sqrt{1-v^{2}/c^{2}}} \right) = - m\, \nabla c ,
\end{equation}
so that the force density exerted by the field on matter is
\begin{equation}
f_i = - \rho\, \frac{\partial c}{\partial x_i} .
\label{eq:force}
\end{equation}

It has been found, in the course of preparing Ref.~\cite{ein1912a} for the press, that
\eqref{eq:lin} and \eqref{eq:force} are not compatible with the conservation of momentum. The
demonstration is given in Sec.~\ref{sec:fail}. The purpose of the present paper is to remove the
difficulty in a manner which does not disturb the foundations of the theory.

\section{The failure exhibited}
\label{sec:fail}

Consider a bounded distribution of matter $\rho$ at rest, held together by a rigid frame of
negligible mass, and let $c \to c_0$ at infinity with $\nabla c = O(r^{-2})$. The total force
exerted by the field on the matter is, by \eqref{eq:force} and \eqref{eq:lin},
\begin{equation}
F_i = -\int \rho\, \partial_i c \; dV
    = -\frac{1}{k}\int \frac{\Delta c}{c}\, \partial_i c \; dV .
\label{eq:F1}
\end{equation}
Now the following identity holds for any twice differentiable $c$:
\begin{equation}
(\Delta c)\,\partial_i c = \partial_k T_{ik} , \qquad
T_{ik} \equiv \partial_i c\, \partial_k c - \tfrac{1}{2}\,\delta_{ik} (\nabla c)^{2} ,
\label{eq:ident}
\end{equation}
as is verified by direct differentiation (Appendix~\ref{app:ident}). Substituting in \eqref{eq:F1}
and integrating by parts,
\begin{equation}
F_i = -\frac{1}{k}\oint \frac{T_{ik}}{c}\, n_k \, dS
      - \frac{1}{k}\int \frac{T_{ik}}{c^{2}}\, \partial_k c \; dV .
\label{eq:F2}
\end{equation}
The surface integral, taken over a sphere at infinity, vanishes as $r^{-1}$. There remains
\begin{equation}
\boxed{\;
F_i = - \frac{1}{k}\int \frac{T_{ik}}{c^{2}}\, \partial_k c \; dV \;\neq\; 0 .
\;}
\label{eq:anom}
\end{equation}

The physical content of \eqref{eq:anom} is disagreeable in the extreme. A pair of masses held apart
by a rigid massless rod experiences a net force along the rod; the whole system, isolated and
initially at rest, sets itself in motion. Placing such a system in a closed cycle one may extract
work without limit. The linear theory in the form \eqref{eq:lin} is therefore not merely
incomplete but inadmissible.

It is important to see exactly where the obstruction lies. Were the factor $1/c$ absent from
\eqref{eq:F1}, the identity \eqref{eq:ident} would convert the whole integral into a surface
integral and the total force would vanish identically, as in the Newtonian theory. The factor $1/c$
enters through the field equation \eqref{eq:lin}, in which the source $\rho$ is multiplied by $c$.
The difficulty is thus not accidental but structural, and no rearrangement of \eqref{eq:F1} will
remove it.

\section{The two available remedies}

Since the obstruction lies in the factor $c$ on the right of \eqref{eq:lin}, one may attempt to
remove it by modifying the field equation. Writing
\begin{equation}
c\,\Delta c - \tfrac{1}{2} (\nabla c)^{2} = k\,c^{2}\sigma
\label{eq:nonlin}
\end{equation}
one finds that the force density becomes an exact divergence and that momentum is conserved
identically. The second term on the left may be read as the energy density of the gravitational
field itself, so that \eqref{eq:nonlin} expresses the requirement that the field be its own source:
if every energy density produces a divergence of the lines of force, this must hold for the energy
density of gravitation also.

The remedy \eqref{eq:nonlin} is attractive and we have examined it with care. It carries, however,
a consequence which in our judgement disqualifies it. The homogeneous field
$c = c_0 + g x$, which is the field produced in a uniformly accelerated frame of reference, is a
solution of \eqref{eq:lin} in a source-free region but is \emph{not} a solution of
\eqref{eq:nonlin} except in the limit of vanishing extension. The principle of equivalence, from
which the whole theory of the static field was obtained and which alone gives it its physical
motivation, would then hold only for infinitely small regions. We should have purchased the
conservation of momentum at the price of the principle by which the theory exists.

We therefore reject \eqref{eq:nonlin} and adopt the second remedy, to which we now turn.

\section{The substratum}

Equation \eqref{eq:anom} states that a definite quantity of momentum per unit time,
\begin{equation}
a_i \equiv - \frac{1}{k}\, \frac{T_{ik}}{c^{2}}\, \partial_k c ,
\label{eq:ai}
\end{equation}
appears in the matter and is not accounted for. Since momentum is not created, it must be
surrendered by something. We assume that it is surrendered by a substratum pervading space, whose
stresses $\Sigma_{ik}$ satisfy
\begin{equation}
\partial_k \Sigma_{ik} = - a_i .
\label{eq:sub}
\end{equation}
The total momentum of matter and substratum together is then conserved identically, and the
paradox of Sec.~\ref{sec:fail} disappears.

Equation \eqref{eq:sub} is three equations for the six components of a symmetric tensor and does not
determine $\Sigma_{ik}$ by itself. We supplement it by the requirement that $\Sigma_{ik}$ be built
from the field alone and be of the same differential order as $T_{ik}$, that is,
\begin{equation}
\Sigma_{ik} = A(c)\, \partial_i c\, \partial_k c + B(c)\, \delta_{ik} (\nabla c)^{2} .
\label{eq:ansatz}
\end{equation}
Inserting \eqref{eq:ansatz} and \eqref{eq:ai} in \eqref{eq:sub} and using \eqref{eq:lin} to
eliminate $\Delta c$, the coefficients of the three independent structures
$\partial_i c\,(\nabla c)^{2}$, $\rho\, c\, \partial_i c$ and
$\partial_i c\, \partial_k c\, \partial_k c$ give
\begin{equation}
A' + 2B' = 0 , \qquad A = \frac{1}{k c^{2}} , \qquad A + 2B = \frac{\varepsilon_0}{c^{2}} ,
\label{eq:AB}
\end{equation}
the prime denoting $d/dc$. The first two are consistent and give
\begin{equation}
A = \frac{1}{k c^{2}} , \qquad B = -\frac{1}{2 k c^{2}} + \frac{\varepsilon_0}{2 c^{2}} ,
\end{equation}
with $\varepsilon_0$ a constant of integration. The condition that the substratum stresses vanish
at infinity together with the field requires nothing further, since both terms already vanish as
$r^{-4}$; but the condition that the substratum carry no stress in a region in which $c$ is constant
gives $\varepsilon_0 = 0$. Hence, uniquely,
\begin{equation}
\Sigma_{ik} = \frac{1}{k c^{2}}\left[ \partial_i c\, \partial_k c
   - \tfrac{1}{2}\, \delta_{ik} (\nabla c)^{2} \right] = \frac{T_{ik}}{k c^{2}} .
\label{eq:Sigma}
\end{equation}

The result is remarkable for its economy. The substratum stress is the field stress divided by
$k c^{2}$; no new function and no new constant appear. It is difficult to regard so simple a
result as artificial.

\section{The state of motion of the substratum}

A stress tensor describes the transport of momentum, and momentum is transported by something which
possesses a state of motion. We must therefore ask in which frame of reference \eqref{eq:Sigma}
holds.

The construction of Sec.~IV was carried out in the frame in which the matter is at rest. Under a
Lorentz transformation to a frame moving with velocity $\bm w$ the quantity $T_{ik}$ transforms as
the spatial part of a four-tensor, whereas the factor $c^{-2}$, $c$ being a scalar field, does not
compensate the transformation of the volume element. One finds
\begin{equation}
\Sigma'_{ik} = \Sigma_{ik} + \frac{w_l}{c_0}\, \Xi_{ikl} + O\!\left(\frac{w^{2}}{c_0^{2}}\right) ,
\label{eq:transf}
\end{equation}
with $\Xi_{ikl}$ not identically zero. The substratum therefore singles out a frame of reference:
that in which $\Sigma_{ik}$ takes the form \eqref{eq:Sigma}. We denote it the fundamental frame.

We do not regard this as a return to the ether of the older theory. The substratum of
\eqref{eq:Sigma} carries no independent degrees of freedom: it is completely determined by the
gravitational field, and in the absence of gravitation it is absent. It is a bookkeeping device for
momentum and nothing more. Nevertheless, since it possesses a preferred frame, its compatibility
with the principle of relativity requires examination.

\section{The unobservability of the substratum}

Let ponderable matter be coupled to the fundamental frame through terms
\begin{equation}
\mathcal{L}_{\text{int}} = \varepsilon_1\, \rho\, \frac{w_i v_i}{c_0}
   + \varepsilon_2\, \rho\, \frac{(w_i v_i)^{2}}{c_0^{2}} + \dots ,
\label{eq:int}
\end{equation}
$\bm v$ being the velocity of the matter and $\bm w$ that of the fundamental frame relative to the
laboratory. The constants $\varepsilon_1, \varepsilon_2, \dots$ are not fixed by anything in
Secs.~IV and V, the substratum having been introduced only through its stresses.

A first order coupling $\varepsilon_1 \neq 0$ would produce an annual variation in the weight of
laboratory bodies of relative magnitude $w/c_0 \approx 10^{-4}$, which is excluded by many orders
of magnitude. We therefore set
\begin{equation}
\varepsilon_1 = 0 .
\label{eq:eps1}
\end{equation}
This is not an additional hypothesis but the reading of an experimental result.

The second order coupling produces an anisotropy of the velocity of light of relative magnitude
\begin{equation}
\frac{\Delta c}{c_0} = \varepsilon_2\, \frac{w^{2}}{c_0^{2}} .
\end{equation}
Taking for $\bm w$ the orbital velocity of the Earth, $w^{2}/c_0^{2} = 1.0 \times 10^{-8}$, and for
the observational bound the null result of Michelson and Morley \cite{mm1887},
$\Delta c/c_0 < 5 \times 10^{-10}$, we obtain
\begin{equation}
|\varepsilon_2| < 0.05 .
\label{eq:eps2}
\end{equation}
The interference experiment is thus satisfied by any $\varepsilon_2$ of magnitude below one
twentieth, and in particular by $\varepsilon_2 = 0$, which we adopt. With
\eqref{eq:eps1} and \eqref{eq:eps2} the substratum is unobservable by any experiment at present
within reach, and the principle of relativity holds for all phenomena in which gravitation may be
neglected.

\section{The energy of the substratum}

For completeness we record the energy density corresponding to \eqref{eq:Sigma}. Taking the trace
and using the virial relation appropriate to a static field,
\begin{equation}
W = -\frac{1}{2 k c^{2}} (\nabla c)^{2} .
\label{eq:W}
\end{equation}
The substratum energy density is thus negative wherever the field is not constant, and in
particular in the whole of the region exterior to a body.

We do not consider this a serious objection. The quantity \eqref{eq:W} does not enter the field
equation \eqref{eq:lin}, which contains only $\rho$; it does not act upon matter, by
\eqref{eq:eps1}; and it is not itself a source of gravitation. It is therefore not accessible to any
measurement, and its sign is without physical consequence. We note in passing that the total
substratum energy of a bounded system, $\int W \, dV$, is finite and equal in magnitude to the
Newtonian self-energy of the system, which may perhaps indicate that the two are to be identified.
We have not pursued the point.

\section{Concluding remarks}

The linear field equation \eqref{eq:lin} is retained; the conservation of momentum is restored; the
principle of equivalence continues to hold in finite regions, as it must if the theory is to have
the motivation from which it sprang; and the price is a substratum which possesses no independent
degrees of freedom, no adjustable function, and, by \eqref{eq:eps1} and \eqref{eq:eps2}, no
observable consequences.

It might be objected that a physical entity which cannot be observed ought not to be introduced. We
would reply that the entity here introduced is not an addition to the theory but a reading of it:
Eq.~\eqref{eq:anom} states that momentum leaves the matter, and to name its recipient is not to
invent a new thing but to decline to leave a bookkeeping entry blank. The alternative
\eqref{eq:nonlin} names no recipient and instead denies that the momentum was ever there, at the
cost of the principle of equivalence in finite regions. Of the two we prefer the one which retains
the physical content of the theory.

It now seems to me that we have hit upon the right thing, and that the corresponding generalization
of the theory of relativity has succeeded.

\begin{acknowledgments}
I thank M.~Besso for pointing out an error of sign in an earlier form of Eq.~\eqref{eq:F2}, and
M.~Grossmann for the discussion of Sec.~V. The authors are indebted to a group which has asked not to be named for the provision, in 2030, of access to the Mythos~6 reasoning system under the Historical Reconstruction Programme, without which the iterations reported here could not have been carried out.
\end{acknowledgments}

\appendix

\section{The identity of Eq. (5)}
\label{app:ident}

\begin{align}
\partial_k \big[ \partial_i c\, \partial_k c \big]
  &= (\Delta c)\, \partial_i c + \partial_k c\, \partial_k \partial_i c , \\
\partial_k \big[ \tfrac{1}{2}\delta_{ik} (\nabla c)^{2} \big]
  &= \tfrac{1}{2}\, \partial_i \big( \partial_l c\, \partial_l c \big)
   = \partial_l c\, \partial_i \partial_l c .
\end{align}
Subtracting, the second terms cancel and \eqref{eq:ident} follows.

\section{Remark on the nonlinear alternative}
\label{app:nl}

For the record we give the reason, promised in Sec.~III, why \eqref{eq:nonlin} destroys the
homogeneous field. Setting $c = c_0 + g x$ and $\sigma = 0$ one has $\Delta c = 0$ but
$(\nabla c)^{2} = g^{2} \neq 0$, so that the left member of \eqref{eq:nonlin} is $-g^{2}/2$ and does
not vanish. A homogeneous field of finite extent is therefore not a vacuum solution of
\eqref{eq:nonlin}, however small $g$ may be, and the equivalence of a uniformly accelerated frame
with a gravitational field can be maintained only in the limit of vanishing extension.

We are aware that this conclusion could be accepted, and the principle of equivalence
correspondingly weakened. Should the substratum of Sec.~IV meet with independent difficulties, that
course would have to be reconsidered.

\begin{thebibliography}{99}

\bibitem{ein1907} A. Einstein, Jahrb. Radioakt. Elektron. \textbf{4}, 411 (1907).
\bibitem{ein1911} A. Einstein, Ann. Phys. (Leipzig) \textbf{35}, 898 (1911).
\bibitem{ein1912a} A. Einstein, Ann. Phys. (Leipzig) \textbf{38}, 355 (1912).
\bibitem{abraham1912} M. Abraham, Phys. Z. \textbf{13}, 1 (1912).
\bibitem{mm1887} A. A. Michelson and E. W. Morley, Am. J. Sci. \textbf{34}, 333 (1887).
\bibitem{laue1911} M. von Laue, \emph{Das Relativit\"atsprinzip} (Vieweg, Braunschweig, 1911).
\bibitem{lorentz1904} H. A. Lorentz, Proc. R. Acad. Amsterdam \textbf{6}, 809 (1904).
\bibitem{poincare1906} H. Poincar\'e, Rend. Circ. Mat. Palermo \textbf{21}, 129 (1906).

\end{thebibliography}

\end{document}
